\documentclass[11pt,a4paper]{article}
\usepackage[utf8]{inputenc}
\usepackage[T1]{fontenc}
\usepackage[spanish,es-minimal,es-noshorthands]{babel}
\usepackage{lmodern}
\usepackage[sfdefault,lining]{FiraSans}
\usepackage[a4paper,top=2.3cm,bottom=2.2cm,left=1.8cm,right=1.8cm,headheight=1.3cm,headsep=0.4cm,footskip=1cm]{geometry}
\usepackage{amsmath,amssymb}
\usepackage{icomma}
\usepackage{graphicx}
\usepackage[table]{xcolor}
\usepackage{tikz}
\usepackage{fontawesome5}
\usepackage{ccicons}
\usepackage{fancyhdr}
\usepackage{lastpage}
\usepackage{tasks}
\usepackage{needspace}
\usepackage[most]{tcolorbox}
\usepackage{array}
\usepackage[hidelinks]{hyperref}
\AtBeginDocument{\hypersetup{pdftitle={Ficha de ejercicios de la unidad 2: Divisibilidad. 1º ESO},pdfsubject={Matemáticas 1º ESO},pdfkeywords={divisibilidad, múltiplos, divisores, criterios de divisibilidad, números primos, factores primos, m.c.m., m.c.d.}}}

\definecolor{navy}{HTML}{1B2748}
\definecolor{accent}{HTML}{E8603A}
\definecolor{teal}{HTML}{1F9E90}
\definecolor{soft}{HTML}{EEF1F8}
\definecolor{muted}{HTML}{5D6685}

\setlength{\parindent}{0pt}
\setlength{\parskip}{3pt}
\settasks{label=\alph*), label-format=\bfseries\color{navy}, label-width=1.3em, item-indent=1.8em, column-sep=0.6em, after-item-skip=2pt, before-skip=3pt, after-skip=4pt}

% ---------- capçalera i peu ----------
\newcommand{\llicencia}{\href{https://creativecommons.org/licenses/by-nc-sa/4.0/deed.es}{\ccbyncsa}\ \ Licencia Creative Commons Atribución-NoComercial-CompartirIgual 4.0 Internacional (CC BY-NC-SA 4.0)}
\pagestyle{fancy}
\fancyhf{}
\fancyhead[L]{\raisebox{-0.35\height}{\includegraphics[height=0.95cm]{logo-gva.png}}}
\fancyhead[C]{\small\color{muted}Ficha de ejercicios · Unidad 2: Divisibilidad · 1\textsuperscript{o} ESO}
\fancyhead[R]{\raisebox{-0.35\height}{\includegraphics[height=0.95cm]{logo-vermellar.png}}}
\fancyfoot[L]{\scriptsize\color{muted}\llicencia}
\fancyfoot[R]{\small\color{muted}\thepage}
\renewcommand{\headrulewidth}{0.4pt}
\renewcommand{\headrule}{\color{muted!40}\hrule height \headrulewidth}
\fancypagestyle{primera}{\fancyhf{}\renewcommand{\headrulewidth}{0pt}
  \fancyfoot[L]{\scriptsize\color{muted}\llicencia}
  \fancyfoot[R]{\small\color{muted}\thepage}}

% ---------- nivells i exercicis ----------
\newcommand{\estrelles}[1]{{\color{accent}\ifcase#1\or\faStar\faStar[regular]\faStar[regular]\or\faStar\faStar\faStar[regular]\or\faStar\faStar\faStar\fi}}
\newcounter{exer}
\newcommand{\ex}[1]{\par\needspace{4\baselineskip}\medskip\stepcounter{exer}\noindent{\color{navy}\bfseries\theexer.}\ \ {\small\estrelles{#1}}\ \ }

\newcommand{\seccio}[2]{\par\needspace{10\baselineskip}\bigskip
  \begin{tcolorbox}[colback=navy,colframe=navy,coltext=white,arc=2mm,boxrule=0pt,left=3mm,right=3mm,top=1.2mm,bottom=1.2mm]
  {\large\bfseries #1\quad #2}
  \end{tcolorbox}\vspace{-2pt}}

\newcolumntype{C}[1]{>{\centering\arraybackslash}p{#1}}

% quadre de recordatori (només a la fitxa d'exercicis)
\newcommand{\recorda}[1]{\par\smallskip\begin{tcolorbox}[colback=accent!7,colframe=accent!55,arc=2mm,boxrule=0.5pt,left=2.5mm,right=2.5mm,top=1mm,bottom=1mm]\small{\color{accent}\faLightbulb[regular]}\ \textbf{Recuerda:}\ #1\end{tcolorbox}}
\newcommand{\res}[1]{#1}

% =====================================================================
% Capçalera FentMates de les fitxes: un color per a cada curs.
% Els colors són els de la pàgina inicial (LEVELS a assets/js/materials.js);
% el logo ve de _eines/logo/pdf/ (python3 _eines/logo/genera.py).
%
% Ús, al final del preàmbul de la unitat:
%   % =====================================================================
% Capçalera FentMates de les fitxes: un color per a cada curs.
% Els colors són els de la pàgina inicial (LEVELS a assets/js/materials.js);
% el logo ve de _eines/logo/pdf/ (python3 _eines/logo/genera.py).
%
% Ús, al final del preàmbul de la unitat:
%   % =====================================================================
% Capçalera FentMates de les fitxes: un color per a cada curs.
% Els colors són els de la pàgina inicial (LEVELS a assets/js/materials.js);
% el logo ve de _eines/logo/pdf/ (python3 _eines/logo/genera.py).
%
% Ús, al final del preàmbul de la unitat:
%   \input{../comu/fentmates.tex}
%   \fmcurs{1eso}
% i, al començament del document:
%   \capcalera{Matemàtiques}{1r ESO}{Fitxa d'exercicis · Unitat 2}{Divisibilitat}
% Opcional abans de \input: \def\fmneed{14} (espai mínim abans de cada secció).
% =====================================================================
\makeatletter
\def\fm@logodir{../../logo/pdf}
\providecommand{\fmneed}{10}
\def\fm@cursnom{}

% clau del curs, color i color del text damunt del color
\newcommand{\fm@curs}[3]{\@namedef{fm@color@#1}{#2}\@namedef{fm@text@#1}{#3}}
\fm@curs{1eso}{E6B84A}{navy}
\fm@curs{2eso}{3B82F6}{white}
\fm@curs{3eso}{FF7A50}{white}
\fm@curs{4eso-a}{A78BFA}{white}
\fm@curs{4eso-b}{6D28D9}{white}
\fm@curs{1batx-mi}{9B1C1C}{white}
\fm@curs{1batx-ccss}{EF6B6B}{white}
\fm@curs{2batx-mii}{1E6B47}{white}
\fm@curs{2batx-ccss}{2BB3A3}{white}

\newcommand{\fmcurs}[1]{%
  \@ifundefined{fm@color@#1}{\PackageError{fentmates}{Curs desconegut: #1}{}}{}%
  \edef\fm@tmp{\noexpand\definecolor{curs}{HTML}{\@nameuse{fm@color@#1}}%
               \noexpand\colorlet{curstext}{\@nameuse{fm@text@#1}}}\fm@tmp
  \colorlet{cursfosc}{curs!72!black}%
  \def\fm@logo{\fm@logodir/vertical-#1.pdf}%
  \def\fm@marca{\fm@logodir/marca-#1.pdf}}

% ---------- capçalera de la primera pàgina ----------
% #1 matèria (petit), #2 curs, #3 tipus de document i unitat, #4 títol
\newcommand{\capcalera}[4]{%
  \gdef\fm@cursnom{#2}%
  \par\noindent
  \begin{tcolorbox}[colback=curs,colframe=curs,boxrule=0pt,arc=2.2mm,width=3.5cm,height=2.15cm,
      valign=center,left=3mm,right=2mm,top=1mm,bottom=1mm,nobeforeafter,box align=center]
    \color{curstext}{\footnotesize #1}\par\vspace{1pt}{\fontsize{19}{21}\selectfont\bfseries #2\par}%
  \end{tcolorbox}\hspace{5mm}%
  \begin{minipage}[c]{\dimexpr\linewidth-3.5cm-5mm-2.9cm\relax}
    {\color{cursfosc}\small\bfseries\MakeUppercase{#3}}\par\vspace{2pt}
    {\color{navy}\fontsize{20}{23}\selectfont\bfseries #4\par}
  \end{minipage}\hfill
  \begin{minipage}[c]{2.6cm}\raggedleft\includegraphics[width=2.6cm]{\fm@logo}\end{minipage}\par
  \vspace{2pt}}

% ---------- barres de secció del color del curs ----------
\renewcommand{\seccio}[2]{\par\needspace{\fmneed\baselineskip}\bigskip
  \begin{tcolorbox}[colback=curs,colframe=curs,coltext=curstext,arc=2mm,boxrule=0pt,left=3mm,right=3mm,top=1.2mm,bottom=1.2mm]
  {\large\bfseries #1\quad #2}
  \end{tcolorbox}\vspace{-2pt}}

% ---------- estrelles de dificultat del color del curs (a Batxillerat no se'n posen) ----------
\@ifundefined{estrelles}{}{%
  \renewcommand{\estrelles}[1]{{\color{curs}\ifcase#1\or\faStar\faStar[regular]\faStar[regular]\or\faStar\faStar\faStar[regular]\or\faStar\faStar\faStar\fi}}}

% ---------- capçalera i peu de les pàgines ----------
% primera pàgina: logos institucionals al peu
\fancypagestyle{primera}{\fancyhf{}\renewcommand{\headrulewidth}{0pt}
  \fancyfoot[L]{\raisebox{-0.32\height}{\includegraphics[height=0.78cm]{logo-gva.png}}\hspace{5mm}%
    \raisebox{-0.32\height}{\includegraphics[height=0.78cm]{logo-vermellar.png}}}
  \fancyfoot[R]{\scriptsize\color{muted}\href{https://creativecommons.org/licenses/by-nc-sa/4.0/deed.ca}{\ccbyncsa}\ \ CC BY-NC-SA 4.0 · fentmates.org\qquad\small\thepage}}
% la resta: el símbol a l'esquerra i el curs a la dreta, amb el títol al centre com abans
\fancyhead[L]{\raisebox{-0.3\height}{\includegraphics[height=0.62cm]{\fm@marca}}}
\fancyhead[R]{\small\bfseries\color{cursfosc}\fm@cursnom}
\renewcommand{\headrule}{\color{curs!45}\hrule height \headrulewidth}
\makeatother

%   \fmcurs{1eso}
% i, al començament del document:
%   \capcalera{Matemàtiques}{1r ESO}{Fitxa d'exercicis · Unitat 2}{Divisibilitat}
% Opcional abans de \input: \def\fmneed{14} (espai mínim abans de cada secció).
% =====================================================================
\makeatletter
\def\fm@logodir{../../logo/pdf}
\providecommand{\fmneed}{10}
\def\fm@cursnom{}

% clau del curs, color i color del text damunt del color
\newcommand{\fm@curs}[3]{\@namedef{fm@color@#1}{#2}\@namedef{fm@text@#1}{#3}}
\fm@curs{1eso}{E6B84A}{navy}
\fm@curs{2eso}{3B82F6}{white}
\fm@curs{3eso}{FF7A50}{white}
\fm@curs{4eso-a}{A78BFA}{white}
\fm@curs{4eso-b}{6D28D9}{white}
\fm@curs{1batx-mi}{9B1C1C}{white}
\fm@curs{1batx-ccss}{EF6B6B}{white}
\fm@curs{2batx-mii}{1E6B47}{white}
\fm@curs{2batx-ccss}{2BB3A3}{white}

\newcommand{\fmcurs}[1]{%
  \@ifundefined{fm@color@#1}{\PackageError{fentmates}{Curs desconegut: #1}{}}{}%
  \edef\fm@tmp{\noexpand\definecolor{curs}{HTML}{\@nameuse{fm@color@#1}}%
               \noexpand\colorlet{curstext}{\@nameuse{fm@text@#1}}}\fm@tmp
  \colorlet{cursfosc}{curs!72!black}%
  \def\fm@logo{\fm@logodir/vertical-#1.pdf}%
  \def\fm@marca{\fm@logodir/marca-#1.pdf}}

% ---------- capçalera de la primera pàgina ----------
% #1 matèria (petit), #2 curs, #3 tipus de document i unitat, #4 títol
\newcommand{\capcalera}[4]{%
  \gdef\fm@cursnom{#2}%
  \par\noindent
  \begin{tcolorbox}[colback=curs,colframe=curs,boxrule=0pt,arc=2.2mm,width=3.5cm,height=2.15cm,
      valign=center,left=3mm,right=2mm,top=1mm,bottom=1mm,nobeforeafter,box align=center]
    \color{curstext}{\footnotesize #1}\par\vspace{1pt}{\fontsize{19}{21}\selectfont\bfseries #2\par}%
  \end{tcolorbox}\hspace{5mm}%
  \begin{minipage}[c]{\dimexpr\linewidth-3.5cm-5mm-2.9cm\relax}
    {\color{cursfosc}\small\bfseries\MakeUppercase{#3}}\par\vspace{2pt}
    {\color{navy}\fontsize{20}{23}\selectfont\bfseries #4\par}
  \end{minipage}\hfill
  \begin{minipage}[c]{2.6cm}\raggedleft\includegraphics[width=2.6cm]{\fm@logo}\end{minipage}\par
  \vspace{2pt}}

% ---------- barres de secció del color del curs ----------
\renewcommand{\seccio}[2]{\par\needspace{\fmneed\baselineskip}\bigskip
  \begin{tcolorbox}[colback=curs,colframe=curs,coltext=curstext,arc=2mm,boxrule=0pt,left=3mm,right=3mm,top=1.2mm,bottom=1.2mm]
  {\large\bfseries #1\quad #2}
  \end{tcolorbox}\vspace{-2pt}}

% ---------- estrelles de dificultat del color del curs (a Batxillerat no se'n posen) ----------
\@ifundefined{estrelles}{}{%
  \renewcommand{\estrelles}[1]{{\color{curs}\ifcase#1\or\faStar\faStar[regular]\faStar[regular]\or\faStar\faStar\faStar[regular]\or\faStar\faStar\faStar\fi}}}

% ---------- capçalera i peu de les pàgines ----------
% primera pàgina: logos institucionals al peu
\fancypagestyle{primera}{\fancyhf{}\renewcommand{\headrulewidth}{0pt}
  \fancyfoot[L]{\raisebox{-0.32\height}{\includegraphics[height=0.78cm]{logo-gva.png}}\hspace{5mm}%
    \raisebox{-0.32\height}{\includegraphics[height=0.78cm]{logo-vermellar.png}}}
  \fancyfoot[R]{\scriptsize\color{muted}\href{https://creativecommons.org/licenses/by-nc-sa/4.0/deed.ca}{\ccbyncsa}\ \ CC BY-NC-SA 4.0 · fentmates.org\qquad\small\thepage}}
% la resta: el símbol a l'esquerra i el curs a la dreta, amb el títol al centre com abans
\fancyhead[L]{\raisebox{-0.3\height}{\includegraphics[height=0.62cm]{\fm@marca}}}
\fancyhead[R]{\small\bfseries\color{cursfosc}\fm@cursnom}
\renewcommand{\headrule}{\color{curs!45}\hrule height \headrulewidth}
\makeatother

%   \fmcurs{1eso}
% i, al començament del document:
%   \capcalera{Matemàtiques}{1r ESO}{Fitxa d'exercicis · Unitat 2}{Divisibilitat}
% Opcional abans de \input: \def\fmneed{14} (espai mínim abans de cada secció).
% =====================================================================
\makeatletter
\def\fm@logodir{../../logo/pdf}
\providecommand{\fmneed}{10}
\def\fm@cursnom{}

% clau del curs, color i color del text damunt del color
\newcommand{\fm@curs}[3]{\@namedef{fm@color@#1}{#2}\@namedef{fm@text@#1}{#3}}
\fm@curs{1eso}{E6B84A}{navy}
\fm@curs{2eso}{3B82F6}{white}
\fm@curs{3eso}{FF7A50}{white}
\fm@curs{4eso-a}{A78BFA}{white}
\fm@curs{4eso-b}{6D28D9}{white}
\fm@curs{1batx-mi}{9B1C1C}{white}
\fm@curs{1batx-ccss}{EF6B6B}{white}
\fm@curs{2batx-mii}{1E6B47}{white}
\fm@curs{2batx-ccss}{2BB3A3}{white}

\newcommand{\fmcurs}[1]{%
  \@ifundefined{fm@color@#1}{\PackageError{fentmates}{Curs desconegut: #1}{}}{}%
  \edef\fm@tmp{\noexpand\definecolor{curs}{HTML}{\@nameuse{fm@color@#1}}%
               \noexpand\colorlet{curstext}{\@nameuse{fm@text@#1}}}\fm@tmp
  \colorlet{cursfosc}{curs!72!black}%
  \def\fm@logo{\fm@logodir/vertical-#1.pdf}%
  \def\fm@marca{\fm@logodir/marca-#1.pdf}}

% ---------- capçalera de la primera pàgina ----------
% #1 matèria (petit), #2 curs, #3 tipus de document i unitat, #4 títol
\newcommand{\capcalera}[4]{%
  \gdef\fm@cursnom{#2}%
  \par\noindent
  \begin{tcolorbox}[colback=curs,colframe=curs,boxrule=0pt,arc=2.2mm,width=3.5cm,height=2.15cm,
      valign=center,left=3mm,right=2mm,top=1mm,bottom=1mm,nobeforeafter,box align=center]
    \color{curstext}{\footnotesize #1}\par\vspace{1pt}{\fontsize{19}{21}\selectfont\bfseries #2\par}%
  \end{tcolorbox}\hspace{5mm}%
  \begin{minipage}[c]{\dimexpr\linewidth-3.5cm-5mm-2.9cm\relax}
    {\color{cursfosc}\small\bfseries\MakeUppercase{#3}}\par\vspace{2pt}
    {\color{navy}\fontsize{20}{23}\selectfont\bfseries #4\par}
  \end{minipage}\hfill
  \begin{minipage}[c]{2.6cm}\raggedleft\includegraphics[width=2.6cm]{\fm@logo}\end{minipage}\par
  \vspace{2pt}}

% ---------- barres de secció del color del curs ----------
\renewcommand{\seccio}[2]{\par\needspace{\fmneed\baselineskip}\bigskip
  \begin{tcolorbox}[colback=curs,colframe=curs,coltext=curstext,arc=2mm,boxrule=0pt,left=3mm,right=3mm,top=1.2mm,bottom=1.2mm]
  {\large\bfseries #1\quad #2}
  \end{tcolorbox}\vspace{-2pt}}

% ---------- estrelles de dificultat del color del curs (a Batxillerat no se'n posen) ----------
\@ifundefined{estrelles}{}{%
  \renewcommand{\estrelles}[1]{{\color{curs}\ifcase#1\or\faStar\faStar[regular]\faStar[regular]\or\faStar\faStar\faStar[regular]\or\faStar\faStar\faStar\fi}}}

% ---------- capçalera i peu de les pàgines ----------
% primera pàgina: logos institucionals al peu
\fancypagestyle{primera}{\fancyhf{}\renewcommand{\headrulewidth}{0pt}
  \fancyfoot[L]{\raisebox{-0.32\height}{\includegraphics[height=0.78cm]{logo-gva.png}}\hspace{5mm}%
    \raisebox{-0.32\height}{\includegraphics[height=0.78cm]{logo-vermellar.png}}}
  \fancyfoot[R]{\scriptsize\color{muted}\href{https://creativecommons.org/licenses/by-nc-sa/4.0/deed.ca}{\ccbyncsa}\ \ CC BY-NC-SA 4.0 · fentmates.org\qquad\small\thepage}}
% la resta: el símbol a l'esquerra i el curs a la dreta, amb el títol al centre com abans
\fancyhead[L]{\raisebox{-0.3\height}{\includegraphics[height=0.62cm]{\fm@marca}}}
\fancyhead[R]{\small\bfseries\color{cursfosc}\fm@cursnom}
\renewcommand{\headrule}{\color{curs!45}\hrule height \headrulewidth}
\makeatother

\fmcurs{1eso}

\begin{document}
\thispagestyle{primera}

\capcalera{Matemáticas}{1\textsuperscript{o} ESO}{Ficha de ejercicios · Unidad 2}{Divisibilidad}
\vspace{10pt}
Nombre y apellidos: \rule{6.6cm}{0.4pt}\hfill Grupo: \rule{1.6cm}{0.4pt}\hfill Fecha: \rule{2.2cm}{0.4pt}
\vspace{8pt}

\begin{tcolorbox}[colback=soft,colframe=soft,arc=2mm,left=3mm,right=3mm,top=1.5mm,bottom=1.5mm]
\small Ejercicios graduados por niveles de dificultad:\quad \estrelles{1}\ básico \qquad \estrelles{2}\ medio \qquad \estrelles{3}\ avanzado.\\[2pt] Resuélvelos en el cuaderno mostrando todos los pasos. La tabla se puede completar directamente en la ficha. Haz los cálculos sin calculadora.
\end{tcolorbox}

% =====================================================
\seccio{1}{La relación de divisibilidad}

\ex{1} Haz la división y di si el primer número es múltiplo del segundo:\begin{tasks}(4)
  \task $84$ y $7$
  \task $95$ y $8$
  \task $132$ y $12$
  \task $200$ y $15$
\end{tasks}\recorda{$a$ es \textbf{múltiplo} de $b$ (y $b$ es \textbf{divisor} de $a$) cuando la división $a:b$ es \textbf{exacta}: resto 0.}

\ex{1} Completa con la palabra \emph{múltiplo} o \emph{divisor}:\begin{tasks}(2)
  \task $35$ es \rule{2cm}{0.4pt} de $7$.
  \task $9$ es \rule{2cm}{0.4pt} de $63$.
  \task $100$ es \rule{2cm}{0.4pt} de $25$.
  \task $13$ es \rule{2cm}{0.4pt} de $52$.
\end{tasks}

\ex{1} A partir de la igualdad $8\cdot9=72$, escribe cuatro frases con las palabras \emph{múltiplo} y \emph{divisor}.

\ex{2} ¿Verdadero o falso? Justifícalo:\begin{tasks}(2)
  \task $48$ es divisible por $6$.
  \task $7$ es divisor de $77$.
  \task $120$ es múltiplo de $9$.
  \task $0$ es múltiplo de $15$.
  \task $1$ es múltiplo de $3$.
\end{tasks}

\ex{2} ¿Cuál es el número más pequeño que hay que sumar a $100$ para que el resultado sea múltiplo de $9$? ¿Y el más pequeño que hay que restarle?

\ex{3} En una división exacta, el divisor es $13$ y el cociente, $17$. ¿Cuál es el dividendo? Con el mismo divisor y cociente, ¿cuál es el dividendo más grande que puede tener una división \emph{no} exacta?

\ex{3} Halla todos los números de dos cifras que son múltiplos de $7$ y acaban en $4$.


% =====================================================
\seccio{2}{Múltiplos de un número}

\ex{1} Escribe los diez primeros múltiplos (sin el 0) de:\begin{tasks}(3)
  \task $7$
  \task $11$
  \task $15$
\end{tasks}

\ex{1} ¿Cuáles de estos números son múltiplos de $6$? \[18\qquad 26\qquad 36\qquad 45\qquad 54\qquad 60\qquad 75\qquad 96\]

\ex{2} Escribe los múltiplos de $13$ comprendidos entre $100$ y $170$.

\ex{2} ¿Cuántos múltiplos de $9$ hay entre $1$ y $200$? ¿Y entre $100$ y $200$?

\ex{2} Escribe los múltiplos comunes de $6$ y de $9$ menores que $100$ (sin el 0). ¿Cuál es el más pequeño?

\ex{3} ¿Cuál es el múltiplo de $17$ más próximo a $300$?

\ex{3} Escribe el múltiplo de $12$ más grande que tenga tres cifras y el más pequeño que tenga cuatro.


% =====================================================
\seccio{3}{Divisores de un número}

\ex{1} Calcula todos los divisores de:\begin{tasks}(4)
  \task $15$
  \task $21$
  \task $28$
  \task $32$
\end{tasks}\recorda{Dividimos entre $1, 2, 3\ldots$ Cada división exacta da dos divisores (el divisor y el cociente). Paramos cuando el cociente es menor o igual que el divisor.}

\ex{1} ¿Cuáles de estos números son divisores de $60$? \[4\qquad 7\qquad 9\qquad 12\qquad 15\qquad 25\qquad 30\]

\ex{2} Calcula todos los divisores de:\begin{tasks}(4)
  \task $48$
  \task $72$
  \task $90$
  \task $100$
\end{tasks}

\ex{2} Escribe los divisores comunes de $24$ y de $36$. ¿Cuál es el más grande?

\ex{2} Pablo tiene 36 cromos y los quiere colocar formando filas, todas con el mismo número de cromos. ¿De cuántas maneras lo puede hacer?

\ex{3} ¿Cuál es el número menor que $50$ que tiene más divisores?

\ex{3} Halla un número que tenga exactamente $4$ divisores y otro que tenga exactamente $5$.


% =====================================================
\seccio{4}{Criterios de divisibilidad}

\ex{1} Completa la tabla: marca con una cruz los números por los que es divisible cada número.\par\smallskip{\renewcommand{\arraystretch}{1.5}\centering\begin{tabular}{|C{2.4cm}|C{1.5cm}|C{1.5cm}|C{1.5cm}|C{1.5cm}|C{1.5cm}|}\hline\rowcolor{soft}\textbf{Número} & $\boldsymbol{\div2}$ & $\boldsymbol{\div3}$ & $\boldsymbol{\div5}$ & $\boldsymbol{\div9}$ & $\boldsymbol{\div11}$\\\hline $3080$ &  &  &  &  & \\\hline $7425$ &  &  &  &  & \\\hline $1188$ &  &  &  &  & \\\hline $2543$ &  &  &  &  & \\\hline $6735$ &  &  &  &  & \\\hline \end{tabular}\par}\recorda{Por \textbf{2}: acaba en cifra par. \ Por \textbf{3}: la suma de las cifras es múltiplo de 3. \ Por \textbf{5}: acaba en 0 o 5. \ Por \textbf{9}: la suma de las cifras es múltiplo de 9. \ Por \textbf{11}: la diferencia entre la suma de las cifras de lugar impar y la de lugar par es 0 o múltiplo de 11.}

\ex{1} Escribe tres números de tres cifras que sean divisibles por $3$ pero no por $9$.

\ex{2} ¿Qué cifras puedes poner en lugar de $\square$ para que $5\square4$ sea divisible por $3$? ¿Y por $9$?

\ex{2} ¿Qué cifra hay que poner en lugar de $\square$ para que $38\square$ sea divisible por $2$ y por $3$ a la vez?

\ex{2} Comprueba, con el criterio, si estos números son divisibles por $11$:\begin{tasks}(4)
  \task $2849$
  \task $9375$
  \task $50\,182$
  \task $13\,575$
\end{tasks}

\ex{3} Halla el número más grande de tres cifras que sea divisible por $2$, por $3$ y por $5$ a la vez.

\ex{3} Halla todas las cifras $a$ y $b$ para que el número $7a2b$ sea divisible por $5$ y por $9$ a la vez.


% =====================================================
\seccio{5}{Números primos y compuestos}

\ex{1} Escribe los divisores de cada número y di si es primo o compuesto:\begin{tasks}(3)
  \task $11$
  \task $15$
  \task $19$
  \task $27$
  \task $31$
  \task $49$
\end{tasks}\recorda{Un número es \textbf{primo} si solo tiene dos divisores: el 1 y él mismo. Si tiene más, es \textbf{compuesto}. El 1 no es ni primo ni compuesto.}

\ex{1} Escribe todos los números primos comprendidos entre $60$ y $100$.

\ex{2} Averigua si son primos dividiéndolos entre $2, 3, 5, 7, 11\ldots$ hasta que el cociente sea menor que el divisor:\begin{tasks}(4)
  \task $87$
  \task $107$
  \task $133$
  \task $149$
\end{tasks}

\ex{2} Escribe cada número como suma de dos números primos:\begin{tasks}(4)
  \task $16$
  \task $24$
  \task $36$
  \task $50$
\end{tasks}

\ex{2} Explica por qué ningún número acabado en $5$, mayor que $5$, puede ser primo. ¿Y por qué ningún número par mayor que $2$ puede ser primo?

\ex{3} ¿El año $2027$ es un número primo? ¿Y el $2028$?

\ex{3} Halla tres números primos consecutivos cuya suma sea también un número primo.


% =====================================================
\seccio{6}{Descomposición en factores primos}

\ex{1} Descompón en factores primos:\begin{tasks}(4)
  \task $36$
  \task $50$
  \task $72$
  \task $81$
\end{tasks}\recorda{Dividimos entre el primo más pequeño posible (2, 3, 5, 7\ldots) hasta llegar a 1, y escribimos los factores repetidos como potencia: $\ 360=2^{3}\cdot3^{2}\cdot5$.}

\ex{1} Descompón en factores primos:\begin{tasks}(4)
  \task $120$
  \task $126$
  \task $196$
  \task $225$
\end{tasks}

\ex{2} Descompón en factores primos:\begin{tasks}(4)
  \task $540$
  \task $784$
  \task $1260$
  \task $2310$
\end{tasks}

\ex{2} ¿Qué número tiene cada descomposición?\begin{tasks}(4)
  \task $2^{4}\cdot3$
  \task $2\cdot3^{2}\cdot7$
  \task $2^{2}\cdot5^{3}$
  \task $3\cdot5\cdot11$
\end{tasks}

\ex{2} Completa las descomposiciones:\begin{tasks}(3)
  \task $2^{\square}\cdot3\cdot5^{\square}=600$
  \task $2^{2}\cdot3^{\square}\cdot7=252$
  \task $2\cdot\square^{2}\cdot5=90$
\end{tasks}

\ex{3} ¿Por qué número, el más pequeño posible, hay que multiplicar $90$ para que el resultado sea un cuadrado perfecto?

\ex{3} Un número tiene la descomposición $2^{2}\cdot3\cdot5$. ¿Qué número es? Escribe todos sus divisores.


% =====================================================
\seccio{7}{Mínimo común múltiplo}

\ex{1} Calcula el m.c.m. escribiendo la lista de múltiplos:\begin{tasks}(4)
  \task $\text{m.c.m.}(3,4)$
  \task $\text{m.c.m.}(6,9)$
  \task $\text{m.c.m.}(8,10)$
  \task $\text{m.c.m.}(5,15)$
\end{tasks}

\ex{1} Calcula el m.c.m. a partir de la descomposición en factores primos:\begin{tasks}(4)
  \task $\text{m.c.m.}(12,20)$
  \task $\text{m.c.m.}(18,30)$
  \task $\text{m.c.m.}(24,36)$
  \task $\text{m.c.m.}(14,21)$
\end{tasks}\recorda{\textbf{m.c.m.}: factores \textbf{comunes y no comunes}, cada uno con el \textbf{mayor} exponente.}

\ex{2} Calcula:\begin{tasks}(4)
  \task $\text{m.c.m.}(40,60)$
  \task $\text{m.c.m.}(36,54)$
  \task $\text{m.c.m.}(45,75)$
  \task $\text{m.c.m.}(28,42)$
\end{tasks}

\ex{2} Calcula:\begin{tasks}(3)
  \task $\text{m.c.m.}(6,8,12)$
  \task $\text{m.c.m.}(10,12,15)$
  \task $\text{m.c.m.}(8,18,27)$
\end{tasks}

\ex{2} Calcula sin descomponer y explica cómo lo has hecho:\begin{tasks}(4)
  \task $\text{m.c.m.}(9,10)$
  \task $\text{m.c.m.}(7,35)$
  \task $\text{m.c.m.}(11,13)$
  \task $\text{m.c.m.}(16,48)$
\end{tasks}

\ex{3} Calcula el m.c.m. de $\ a=2^{3}\cdot3^{2}\cdot7\ $ y $\ b=2^{2}\cdot3^{3}\cdot5$.

\ex{3} El m.c.m. de dos números es $60$ y uno de ellos es $12$. ¿Qué valores puede tener el otro?


% =====================================================
\seccio{8}{Problemas de m.c.m.}

\ex{1} Un faro se enciende cada $9$ segundos y otro, cada $12$. Si ahora se han encendido los dos a la vez, ¿dentro de cuánto tiempo volverán a coincidir?\recorda{Si hay cosas que se \textbf{repiten} y queremos saber cuándo volverán a \textbf{coincidir}, o buscamos el número \textbf{más pequeño} que es múltiplo de varios, usamos el \textbf{m.c.m.}}

\ex{1} Lucía entrena a fútbol cada $3$ días y va a natación cada $4$. Hoy ha hecho las dos cosas. ¿Dentro de cuántos días las volverá a hacer el mismo día?

\ex{2} Tres amigos van al cine: uno cada $6$ días, otro cada $8$ y el otro cada $10$. Si hoy han coincidido, ¿dentro de cuántos días volverán a coincidir?

\ex{2} ¿Cuál es el número más pequeño de lápices que se puede empaquetar en cajas de $12$, de $15$ o de $18$ sin que sobre ninguno?

\ex{2} Dos autobuses salen juntos de la parada a las 7:30. Uno vuelve cada $20$ minutos y el otro, cada $25$. ¿A qué hora volverán a coincidir?

\ex{3} Si agrupamos las galletas de una caja de $6$ en $6$, de $8$ en $8$ o de $10$ en $10$, siempre sobran $3$. ¿Cuántas galletas hay como mínimo? ¿Y si sabemos que hay entre $200$ y $300$?

\ex{3} Una rueda dentada de $15$ dientes engrana con otra de $24$ dientes. ¿Cuántas vueltas tiene que dar cada una para que vuelvan a estar en la posición inicial?


% =====================================================
\seccio{9}{Máximo común divisor}

\ex{1} Calcula el m.c.d. escribiendo la lista de divisores:\begin{tasks}(4)
  \task $\text{m.c.d.}(8,12)$
  \task $\text{m.c.d.}(15,20)$
  \task $\text{m.c.d.}(18,27)$
  \task $\text{m.c.d.}(14,25)$
\end{tasks}

\ex{1} Calcula el m.c.d. a partir de la descomposición en factores primos:\begin{tasks}(4)
  \task $\text{m.c.d.}(24,36)$
  \task $\text{m.c.d.}(30,45)$
  \task $\text{m.c.d.}(42,56)$
  \task $\text{m.c.d.}(60,90)$
\end{tasks}\recorda{\textbf{m.c.d.}: \textbf{solo} los factores \textbf{comunes}, cada uno con el \textbf{menor} exponente. Si no hay ninguno común, el m.c.d. es 1 (\emph{primos entre sí}).}

\ex{2} Calcula:\begin{tasks}(4)
  \task $\text{m.c.d.}(84,126)$
  \task $\text{m.c.d.}(120,180)$
  \task $\text{m.c.d.}(96,144)$
  \task $\text{m.c.d.}(105,165)$
\end{tasks}

\ex{2} Calcula:\begin{tasks}(3)
  \task $\text{m.c.d.}(12,18,30)$
  \task $\text{m.c.d.}(24,40,56)$
  \task $\text{m.c.d.}(45,75,105)$
\end{tasks}

\ex{2} ¿Cuáles de estas parejas son primos entre sí?\begin{tasks}(4)
  \task $8$ y $21$
  \task $15$ y $25$
  \task $14$ y $27$
  \task $22$ y $33$
\end{tasks}

\ex{3} Calcula el m.c.d. y el m.c.m. de $36$ y $60$, y comprueba que $\text{m.c.d.}\cdot\text{m.c.m.}=36\cdot60$.

\ex{3} El m.c.d. de dos números es $8$ y su m.c.m., $96$. Si uno de los números es $24$, ¿cuál es el otro?


% =====================================================
\seccio{10}{Problemas de m.c.d.}

\ex{1} Tenemos $30$ manzanas y $45$ naranjas y queremos hacer el máximo número de lotes iguales, sin que sobre ninguna fruta. ¿Cuántos lotes podemos hacer? ¿Qué llevará cada lote?\recorda{Si queremos \textbf{repartir} o \textbf{cortar} en partes iguales, \textbf{sin que sobre nada}, y las partes tienen que ser \textbf{lo más grandes posible} (o el máximo número de grupos), usamos el \textbf{m.c.d.}}

\ex{1} Queremos cortar dos cuerdas de $36$ m y $48$ m en trozos iguales, lo más largos posible. ¿Cuánto medirá cada trozo? ¿Cuántos trozos tendremos?

\ex{2} Queremos embaldosar un suelo de $360$ cm por $280$ cm con baldosas cuadradas, lo más grandes posible, sin cortar ninguna. ¿Cuánto medirá el lado de cada baldosa? ¿Cuántas baldosas hacen falta?

\ex{2} Una librería tiene $120$ libros de cuentos, $84$ de poesía y $60$ cómics. Quiere hacer paquetes con el mismo número de libros, todos de un solo tipo y lo más grandes posible. ¿Cuántos libros habrá en cada paquete? ¿Cuántos paquetes hará?

\ex{2} Di si hay que calcular el m.c.m. o el m.c.d. y resuelve:\begin{tasks}(1)
  \task Un reloj suena cada $15$ minutos y otro, cada $25$. Si ahora han sonado juntos, ¿cuándo volverán a coincidir?
  \task Queremos repartir $48$ bolígrafos y $64$ lápices en el máximo número de estuches iguales, sin que sobre ninguno.
\end{tasks}

\ex{3} Una maestra reparte $75$ cromos y $105$ pegatinas entre el alumnado, de manera que todos reciben la misma cantidad de cada cosa. Le sobran $3$ cromos y $9$ pegatinas. ¿Cuántos alumnos puede haber, como máximo? ¿Cuántos cromos y pegatinas recibe cada uno?

\ex{3} Queremos plantar árboles alrededor de un jardín rectangular de $72$ m por $48$ m, a la misma distancia entre ellos, la mayor posible y con un árbol en cada esquina. ¿A qué distancia los plantaremos? ¿Cuántos árboles hacen falta?

\vfill
\begin{center}\color{muted}\small\faCheckCircle[regular]\ \ ¡Ánimo! Si te cuesta algún ejercicio, repasa los apuntes interactivos de la unidad.\end{center}

\end{document}
